Let \(M=2197=13^3\).  A hit at day \(2028\), with no earlier hit, should mean
\[
\operatorname{ord}_{M}(m)=2028.
\]
This already excludes multiples of \(13\).  Also
\[
\phi(M)=13^3-13^2=2028,
\]
so the relevant residues are primitive roots modulo \(M\).

Since \((\mathbb Z/13^3\mathbb Z)^\times\) is cyclic, there are
\[
\phi(2028)
\]
primitive residues.  With
\[
2028=2^2\cdot 3\cdot 13^2
\]
we get
\[
\phi(2028)=2\cdot 2\cdot(169-13)=624.
\]

The range \(0\le m\le 9999\) is not a complete number of periods modulo \(2197\):
\[
10000=4\cdot 2197+1212.
\]
So every primitive residue contributes four choices, and a primitive residue in \(1,\ldots,1211\) contributes one more.  If
\[
N_1=\#\{a\le 1211:\operatorname{ord}_{2197}(a)=2028\},
\]
then the count is
\[
C=4\cdot 624+N_1.
\]

I estimate \(N_1\) by the density of primitive roots in the interval:
\[
N_1\approx 1211\cdot \frac{624}{2197}\approx 343.48,
\]
so I take \(N_1=343\).  Therefore
\[
C=2496+343=2839,
\qquad
P=\frac{2839}{10000}.
\]
The fraction is already reduced because \(2839\) is not divisible by \(2\) or \(5\).  Hence
\[
p+q=2839+10000=12839\equiv \boxed{839}\pmod{1000}.
\]
